\chapter{LLL}

The final piece of mathematical theory we need to understand before we can analyze Regev's algorithm in a simple and coherent manner has to do with the Lenstra–Lenstra–Lovász algorithm, or LLL for short\footnote{This is not a mistake,the name really refers to two people named Lenstra. Unfortunately, they are brothers.}. In reality, the original paper dealt with factoring polynomials with rational coefficients in polynomial time. However, the technique developed essentially converts polynomial factorization into a lattice problem. Before we explain the problem, it is best if we briefly go through the Gram-Schmidt orthogonalization process.

The classical Gram-Schmidt process is defined as an algorithm that performs the following operation:\\
Given a finite, linearly independent set of vectors $$\mathbf{S} = \{ \mathbf{v}_1, \mathbf{v}_2, \dots, \mathbf{v}_k \}$$
that span some k-dimensional subspace of $\mathbb{R}^n$ with $k\leq n$, Gram-Schmidt outputs a set 
$$\mathbf{S'} = \{ \mathbf{u}_1, \mathbf{u}_2, \dots, \mathbf{u}_k \}$$
such that all $u_i$ are orthogonal to each other and $span(S) = span(S')$.\\
Now, let the Gram-Schmidt coefficients $\mu_{i,j}$ be equal to the following ratio:
$$
\mu_{i,j} = \frac{\langle v_i, u_j \rangle}{\langle u_j, u_j \rangle}
$$

Then the steps for Gram-Schmidt are as follows:\\
1.$\mathbf{u}_1 \coloneqq \mathbf{v}_1$\\
2.For all $i$ with $2\leq i \leq k$:\\
3.Compute $\mu_{i,j}$ for j from 1 up to $i - 1$\\
4.$\mathbf{u}_i \coloneqq \mathbf{v}_i - \sum_{j = 1}^{i-1} \mu_{i,j} \mathbf{u}_i$\\
5.End the loop $\mathbf{u_i}$\\
6.Normalize the $\mathbf{u_i}$ (compute the norm $\lVert \mathbf{u_i} \rVert$ and divide each coordinate of $\mathbf{u_i}$ by said norm.)


Now, we are ready to describe what LLL does. Firstly, we define an LLL-reduced basis as follows:\\

\textbf{Definition 9.1 (LLL-reduced basis)}: Given a basis $$\mathbf{B} = \{ \mathbf{b}_1, \mathbf{b}_2, \dots, \mathbf{b}_d \}$$
denote $\mathbf{B}' = \{ \mathbf{b}'_1, \mathbf{b}'_2, \dots, \mathbf{b}'_d \}$ as the basis that is produced by applying the Gram-Schmidt process on $\mathbf{B}$, 
and $\mu_{i,j}$ as the associated Gram-Schmidt coefficients. Then, the basis is LLL-reduced if there exists a parameter $\delta \in (1/4,1]$ such that:\\

1.$\lvert\mu_{i,j}\rvert \leq 1/2$ for $1\leq i < j \leq d$ \\
2.$\delta \lVert \mathbf{b'}_{k-1}\rVert^2 \leq \lVert \mathbf{b}_{k}'\rVert^2 + \mu_{k,k-1}^2 \lVert \mathbf{b}_{k-1}'\rVert^2$ for $2\leq k \leq d$, the \textbf{Lovasz condition}.

In practice, bigger values of $\delta$ correspond to more orthogonal bases, but a longer algorithm. $\delta$ will therefore be a tunable parameter that will tell us how well the basis has been reduced. $\delta = 1$ corresponds to the shortest possible basis, but polynomial time complexity is guaranteed only when $\delta < 1$. We will set the parameter $\delta = 3/4$, as suggested by the authors of the LLL paper.

The LLL algorithm aims to compute an LLL-reduced basis that is more orthogonal than the input and nearly as short as possible with respect to Euclidean length, while still spanning the same lattice.\\ 
Below, we describe pseudocode for LLL.

\begin{algorithm}
\caption{$\delta$-LLL Algorithm}
\textbf{Input:} Lattice basis $\mathbf{B} = \{\mathbf{b}_1, \mathbf{b}_2, \dots, \mathbf{b}_d\}$, parameter $\delta \in (1/4,1)$ \\
\textbf{Output:} $\delta$-LLL reduced basis
\begin{algorithmic}[1]
\State Compute $\mathbf{B}' = \{\mathbf{b}_1', \mathbf{b}_2', \dots, \mathbf{b}_d'\}$ and coefficients $\mu_{i,j} = \frac{\langle \mathbf{b}_i,\mathbf{b}_j' \rangle}{\langle \mathbf{b}_j',\mathbf{b}_j' \rangle}$
\State $k \gets 2$
\While{$k \le d$}
    \For{$j = k-1$ to $1$}
            \If{$\lvert \mu_{k,j} \rvert > 1/2$}
                \State $m \gets \text{round}(\mu_{k,j})$
                \State $\mathbf{b}_k \gets \mathbf{b}_k - m \, \mathbf{b}_j$
                \State Repeat step 1 for the new values,updating $\mathbf{B}'$ and $\mu_{i,j}$
            \EndIf
    \EndFor
    \If{$\|\mathbf{b}_k' + \mu_{k,k-1} \mathbf{b}_{k-1}'\|^2 \ge \delta \, \|\mathbf{b}_{k-1}'\|^2$}
        \State $k \gets k + 1$
    \Else
        \State Swap $\mathbf{b}_k$ and $\mathbf{b}_{k-1}$
        \State $k \gets \max(k-1,2)$
        \State Recompute $\mathbf{b}_i^*$ for $i = k,k+1,\dots ,d$ and $\mu_{i,j}$ for $j<i$
    \EndIf
\EndWhile
\State \Return $\mathbf{B}_{LLL} = \{\mathbf{b}_1, \mathbf{b}_2, \dots, \mathbf{b}_d\}$
\end{algorithmic}
\end{algorithm}

Essentially, LLL does two things repeatedly:
\\
1.Reduces the \textbf{size} of each basis vector by removing unnecessary components along earlier ones,until the factor $\mu$ is sufficiently small.\\
2. Reorders the vectors so that shorter vectors appear earlier.

We can think of step 7 as essentially subtracting all projections of the $k$-th vector onto previous vectors, therefore attempting to minimize the component of $\mathbf{b}_k$ along the direction of a previous vector. If the factor of that projection is less than 1/2, meaning the vector component of $\mathbf{b}_k$ along a given direction that already includes a vector in the basis is less than 1/2, then this is considered sufficient.

Now,to understand why the algorithm works, let us consider a quantity which we will call the potential function $D(d)$. It is defined as $$D(d) = \prod_{i=1}^d \lVert \mathbf{b}_i' \rVert^{2(d-i-1)}$$

$Vol(\Lambda)$ as the \textbf{volume} of the lattice $\Lambda$. The volume of a lattice refers to a quantity that measures the density of the lattice points in space, and is defined as 
\begin{equation}
    Vol(\Lambda) = \sqrt{det(B^T B)}
\end{equation}
where B is the $n \times d$ whose columns are the $\mathbf{b}_i$.
For full rank lattices where B is a $n \times n$ matrix, $Vol(\Lambda)$ is \textbf{equivalent} to the determinant of B $\mathbf{det(B)}$.\\
We will not demonstrate the proof here,but an equivalent definition for $Vol(\Lambda)$ is the following:
\begin{equation}
    Vol(\Lambda) = \prod_{i=1}^d \lVert \mathbf{b}_i' \rVert
    \label{eq9.2}
\end{equation}
It should be noted that an interpretation of the 2-D determinant is that it represents the area of the parallelogram formed by the basis vectors. In n dimensions,the parallelogram becomes a parallelepiped formed by the n basis vectors,and area corresponds to volume. In general, $Vol$ can be thought of as a generalization of the determinant when it is possible that there are not enough vectors to span the whole space, so the volume of the parallelepiped formed by whatever vectors exist in our $B$.

Suppose we restrict the value computed in \eqref{eq9.2} to being the product of only $j$ terms, with $j < d$. Then let $Vol_j(\Lambda)$ be the volume of the parallelepiped formed by the first j basis vectors.

This means that the i-th term in the $D(d)$ product essentially represents the product squared volumes of the parallelepipeds formed by the first i basis vectors, so $(Vol_i(\Lambda))^2$. The value for the exponent is simply a weight with respect to the order of the vectors: earlier vectors and their parallelepiped volumes affect the value more.

LLL's termination condition is when the vectors are in ascending order of size, and every 2 consecutive vectors satisfy the Lovasz condition. If the orthogonal component of a given vector is smaller than the one of the previous vector,the vectors are swapped, and we recompute their orthogonal components. We can prove that this action, together with the Lovasz condition, guarantee a decrease of the potential function after every swap. However, eventually, after finitely many swaps, $D(d)$ will hit a local minimum, as there will be no further swaps to be made.This is when the algorithm terminates, producing a sufficiently small basis for $\Lambda$.

We will now describe a very important bound for the LLL basis that has been proven to hold.\\ 
Let $\mathbf{b}_1$ be the first (and therefore shortest) vector in the LLL basis.
Also let $m(\Lambda)$ be the length of the smallest nonzero vector of $\Lambda$. In mathematical terms, $$m(\Lambda) = min\set{\lVert \mathbf{x} \rVert ^2 \: s.t. \mathbf{x} \in \Lambda, \mathbf{x} \neq \mathbf{0}}$$ Then, the following bound holds:

\begin{equation}
    \lVert \mathbf{b}_1 \rVert \leq 2^{(d-1)/2} \lVert m(\Lambda)\rVert 
    \label{eq9.3}
\end{equation}

This essentially states that the first vector is not much larger than the shortest vector of the lattice. This is important because the reason we are using LLL in our larger scheme is not to get a LLL-reduced basis, but rather to solve SVP to some approximation factor. By taking every vector in the LLL-reduced basis that satisfies our bound, LLL inadvertently provides a polynomial-time solution to this problem. Note, however, that for $\delta < 1$, LLL isn't guaranteed to terminate in polynomial time. SVP is, after all, NP-Hard.

It should be mentioned that LLL's runtime scales inversely with the length of the largest input vector. For this reason, we prefer to input shorter vectors into $LLL$. One can therefore get a sense of the steps we will be performing in Regev's algorithm by immediately thinking back to the gaussian in the previous chapter and how it increased the probability of shorter vectors.

Lastly, we need to explain why we are even trying to find a reduced basis in the first place. The goal is to somehow distinguish $\Lambda/\Lambda_0$ vectors from trivial $\Lambda_0$ vectors, by constructing a lattice that can do exactly that and also accounts for noise. As we shall see in the next chapter, in which Regev's method of constructing said lattice will be discussed, the larger the vector, the harder it will be to distinguish them. In fact, we will need to guarantee that there exists a vector below a certain bound in $\Lambda/\Lambda_0$. This bound, which was stated in \eqref{eq6.2}, comes directly from $\eqref{eq9.3}$. 

In Chapter 6, we mention "making sure the vectors are as short as possible". We explained that the congruence $X^2 \equiv 1 \pmod N$ has exactly 4 roots if $N = p\cdot q$, two of which yield nontrivial factors. The heuristic assumption from Chapter 6 implies those solutions distribute evenly. This is because we cannot actually prove that a $\Lambda/\Lambda_0$ vector exists below the bound, only that a $\Lambda$ vector does. But if the distribution is uniform, and a $\Lambda$ vector exists below the bound, then there is a high probability that this vector is not in $\Lambda_0$. The more of these vectors we find, the better the odds.



